\documentclass{ltjsbook}
\usepackage{docmute} 
\usepackage{../../myStyle} 

\begin{document}
\chapter{因子分析法}
\label{L08_FacorAnalysis}


すべての正方行列に，こういった「変換」という解釈ができるのであれば，相関行列にも同様のことがいえるでしょう。相関行列は正方行列ですので，固有値分解できます。相関行列を固有値分解することは，相関行列の中に潜む次元(dimension)を抽出してくることです。固有ベクトル(因子負荷行列)は，正方行列によって変換される，変換先の単位ベクトルのことだったのです。そして，固有値はその次元のゆがみ(重み，重要性)という意味があったのです。


分散共分散行列のトレースは，分散の総和を意味します。項目同士の関係を表した行列であれば，分散はその項目から得られる情報の大きさであり，それを総和するということは，その調査研究・項目群から得られる情報の総和であると言ってもいいでしょう。
相関行列のトレースは，対角項に入っているのが$r_{ii}=1.0$ですから，項目の数と一致します。1つの項目の情報量を1.0に基準化してN項目分の情報がある，ということを表しています。


固有値と行列の関係は冒頭で示したように，$\bm{Ax}=\lambda \bm{x}$であり，正方行列の特徴をスカラーにしてしまうというものです。得られる$N$個の固有値は，元の正方行列のエッセンスをスカラーにして表現しているわけです。ここで固有値の大きさの順に並べ替えて，順に$\lambda_1,\lambda_2 \cdots \lambda_n$としたとしましょう。ここで固有ベクトルの大きさを$\bm{xx'}=1$となるように選んでやれば，元の行列を次のように書き換えることができます。

\[
    \bm{Axx'} = \lambda \bm{xx'}
\]

ここから，次のように行列$\bm{A}$を分解して行くことができます。

\[
    \bm{A}= \lambda_1 \bm{x}_1 \bm{x}_1' + \lambda_2 \bm{x}_2 \bm{x}_2' + \cdots \lambda_m \bm{x}_m \bm{x}_m'= \sum \lambda_{i} \bm{x}_{i} \bm{x}_i'
\]

これらをまとめると，固有値は元の行列の対角要素の情報を表したものであり，全体の情報量はそのままに重要度の順に並べ替えたものである，といえるでしょう。
これこそ\keytermL{因子分析}{いんしぶんせき}で取り出そうとしている因子であり，固有値の大きさはその因子の重要度である，ということがわかります！
\end{document}